To validate the reliability and operational success of the Deltx platform, our evaluation strategy will be focusing on the three main layers of our system.

\section{AI Code Detection Evaluation}

Since our AI detection module relies on likelihood and log-probability-based approaches, we will evaluate its performance against a baseline dataset containing known mixtures of human-written and LLM-generated Python code. We will quantitatively measure its accuracy using standard classification metrics: Precision, Recall, and the F1-Score. Our target is to achieve a baseline that aligns with modern state-of-the-art probabilistic detectors.

\section{Predictive Modeling Accuracy}

The core forecasting capability of our PatchTST engine will be evaluated using standard statistical time-series metrics. We will utilize \textbf{Mean Absolute Error (MAE)} to measure the average deviation of our predicted 0-100 quality scores from the actual scores. Additionally, we will calculate the \textbf{Root Mean Square Error (RMSE)} to penalize larger, catastrophic forecasting errors—such as predicting a highly stable codebase when a major quality drop can be observed.

\section{Explainability and Practical Usefulness}

Because our predictive model operates strictly within a metric space—analyzing time-series trends of quality scores, AI percentages, and codebase churn—our interpretation layer will focus specifically on temporal and metric feature attribution using SHAP.

Specifically, we evaluate whether the explainability layer can identify the most influential historical time windows and associated metric shifts that contributed to the forecasted decline.

For example, a predicted drop in Maintainability may be attributed to a combination of increased cognitive complexity and elevated AI-generated code proportion occurring within a specific commit window, which the model identifies as having the strongest influence on the forecast. 

For the evaluation of this layer, we will be using \textbf{human expert validation}, where experienced software engineers will assess the relevance and accuracy of the SHAP explanations in relation to the actual code changes and quality metrics observed in the historical data.